% [R19] Source-attributed presentation derivative of the bilingual V4.2.14 handbooks.
% Presentation layer: Chinese examples, span boundaries, labels, and execution versions remain unchanged.
% G1--G15 are display identifiers, not new normative rule IDs or a relabelled V4.2.9 image.
% Crosswalk and source hashes: handbooks/ALIGNMENT_R18.md. No annotations are changed.
\definecolor{guideL}{HTML}{DED9F7}
\definecolor{guideK}{HTML}{C6ECEE}
\definecolor{guideS}{HTML}{FFF0B0}
\definecolor{guideT}{HTML}{CDEACF}
\newcommand{\gspan}[2]{{\setlength{\fboxsep}{1.4pt}\colorbox{guide#1}{\strut [#2]\,\textsf{[#1]}}}}
\newcommand{\gcontext}[1]{\textcolor{black!75}{#1}}
\newcommand{\gsource}[1]{\,\textsuperscript{\textsf{[#1]}}}
\newcommand{\ggloss}[1]{\par\emph{#1}}

\section{Appendix A. Coding Guide and Examples}
\phantomsection
\label{app:a}
\label{sec:handbook-current}
This appendix summarizes the v4.2.14 rule series with bilingual examples. Figure~\ref{fig:illustrative-annotations} shows the preserved human-reference annotations. The full \href{https://github.com/AlfredJamesLi/chinese-skillspan-benchmark/blob/main/notes/handbooks/reader_20261007/README.md}{Chinese and English handbooks, dated 7 October 2026}, provide further examples and identify cases awaiting review. Section~A.4 distinguishes these reading materials from the versions associated with the reported studies.

\subsection{A.1 Label definitions and notation}
Table~\ref{tab:guide-types} defines the four labels with examples. Square brackets mark selected source substrings, followed by their type codes. Brackets, colours, and type codes are display aids rather than source characters; brackets keep boundaries identifiable in grayscale. The examples illustrate annotation decisions and are not an additional evaluation sample.

\noindent\textit{Source key.} E: ESCO's \href{https://esco.ec.europa.eu/en/about-esco/escopedia/escopedia/skills-pillar}{skills hierarchy} and \href{https://esco.ec.europa.eu/en/about-esco/escopedia/escopedia/transversal-knowledge-skills-and-competences}{transversal concepts} \citep{ESCOv1_2_2024}; Q: EQF definitions \citep{council2017eqf}; Z: related annotation guidance in SkillSpan, Appendix B \citep{zhang2022skillspan}; I: unitization and categorization in agreement analysis \citep{artstein2008agreement}; P: our V4.2.14 operational decision. Combined marks indicate adaptation, not verbatim rules; Chinese examples and exact boundaries are project illustrations.

\Needspace{8\baselineskip}
\subsection{A.2 General annotation decisions}
G1--G15 are quick-reference identifiers. Their source sections and amendment IDs (R08--R23) are mapped in the release; they do not replace the handbook's rule numbering.

\begingroup
\small\setlength{\tabcolsep}{4pt}\renewcommand{\arraystretch}{1.13}
\begin{longtable}{@{}P{0.07\linewidth}P{0.49\linewidth}P{\dimexpr0.44\linewidth-16pt\relax}@{}}
\caption{Annotation principles for competency mentions and span boundaries.}\label{tab:guide-boundaries}\\
\toprule\textbf{Rule} & \textbf{Annotation principle} & \textbf{Example and English gloss}\\\midrule
\endfirsthead
\multicolumn{3}{@{}l}{\textit{Table \thetable\ (continued)}}\\
\toprule\textbf{Rule} & \textbf{Annotation principle} & \textbf{Example and English gloss}\\\midrule
\endhead


G1 & The annotation scheme distinguishes L, K, S, and T. Collapsed categories are used only in secondary evaluation.\gsource{E,P} & \gspan{L}{英语六级}\ggloss{College English Test Band 6}\par\smallskip\gspan{K}{计算机专业}\ggloss{Computer science major}\par\smallskip\gspan{S}{数据分析}\ggloss{Data analysis}\par\smallskip\gspan{T}{责任心}\ggloss{Sense of responsibility} \\[5pt]
G2 & Final annotations contain non-overlapping spans. A sentence may contain several independent mentions, with each character assigned to at most one span.\gsource{P} & \gcontext{熟悉使用}\gspan{S}{Word}\gcontext{，}\gspan{S}{Excel}\gcontext{，}\gspan{S}{PPT}\gcontext{等办公软件}\ggloss{Familiar with using Word, Excel, PowerPoint, and other office software}\par\smallskip No additional covering list span. \\[5pt]
G3 & Each mention is a contiguous source substring that preserves spelling and internal spaces. Any normalization is stored separately.\gsource{Z,P} & \gcontext{熟悉 }\gspan{S}{Jekins}\ggloss{Familiar with Jekins}\par\smallskip Keep the source spelling, not \texttt{Jenkins}. \\[5pt]
G4 & Annotators select the shortest span that expresses a complete requirement. Words and identifiers remain intact; semantic completeness takes priority over length, with no fixed length limit.\gsource{Z,P} & \gspan{S}{支持服务}\ggloss{Support services}\par\smallskip The truncated \zh{支持服} is incomplete. \\[5pt]
G5 & Generic proficiency triggers usually remain outside the span, while heads needed to express the requirement remain inside. Excluded triggers can still inform type assignment.\gsource{Z,P} & \gcontext{会使用}\gspan{S}{Excel}\ggloss{Able to use Excel}\par\smallskip\gspan{T}{会聊天}\ggloss{Can hold a conversation}\par\smallskip\gcontext{有}\gspan{T}{服务意识}\ggloss{Has a service-oriented mindset} \\[5pt]
G6 & Coordinated items are separated when each remains meaningful in context. Shared qualifiers may stay outside the spans, and implied words are not added.\gsource{Z,P} & \gcontext{熟悉使用}\gspan{S}{Word}\gcontext{，}\gspan{S}{Excel}\gcontext{，}\gspan{S}{PPT}\gcontext{等办公软件}\ggloss{Familiar with using Word, Excel, PowerPoint, and other office software}\par\smallskip Shared-experience case: see A.3. \\[5pt]
G7 & A necessary action, object, or shared head remains within the span when splitting would change its meaning. Span length alone does not determine type.\gsource{Z,P} & \gspan{S}{优化曝光与转化率}\ggloss{Optimize exposure and the conversion rate}\par\smallskip Isolated metric names would lose the optimization activity. \\
\bottomrule
\end{longtable}
\endgroup

\begingroup
\small\setlength{\tabcolsep}{4pt}\renewcommand{\arraystretch}{1.13}
\begin{longtable}{@{}P{0.07\linewidth}P{0.49\linewidth}P{\dimexpr0.44\linewidth-16pt\relax}@{}}
\caption{Annotation principles for scope, experience, type, and unresolved records.}\label{tab:guide-scope}\\
\toprule\textbf{Rule} & \textbf{Annotation principle} & \textbf{Example and English gloss}\\\midrule
\endfirsthead
\multicolumn{3}{@{}l}{\textit{Table \thetable\ (continued)}}\\
\toprule\textbf{Rule} & \textbf{Annotation principle} & \textbf{Example and English gloss}\\\midrule
\endhead


G8 & Benefits, company promotion, administrative instructions, experience durations, and pure outcomes are outside competency scope on their own. Mixed clauses are assessed separately.\gsource{Z,P} & \zh{五险一金，带薪年假} $\rightarrow$ \texttt{[]};\ggloss{Five types of social insurance and a housing fund; paid annual leave}\par\smallskip\gspan{S}{分析推广效果}\gcontext{，保证流量/用户增长}\ggloss{Analyze the effectiveness of promotional activities; ensure traffic/user growth} \\[5pt]
G9 & Experience wording is omitted only when the remaining activity preserves the requirement. Necessary practical-experience wording is retained. Explicit industry background follows the handbook's broad-K convention. Experience wording is never added to the source.\gsource{P} & \gspan{S}{机器人竞赛经验}\ggloss{Experience in robotics competitions}\par\smallskip Retain the suffix if a bare competition name changes the requirement; do not label generic \zh{经验} (experience) alone. \\[5pt]
G10 & Educational qualifications include their threshold wording. Language examinations and certificates receive L; non-language qualifications receive K.\gsource{Z,P} & \gspan{K}{本科及以上学历}\ggloss{Bachelor's degree or higher}\par\smallskip\gspan{L}{大学英语6级}\ggloss{College English Test Band 6}\par\smallskip\gcontext{了解}\gspan{K}{ISO 27001}\ggloss{Familiar with ISO 27001}\par\smallskip This denotes familiarity with the standard rather than personal certification. \\[5pt]
G11 & Tool and method mentions are typed from their predicate context. Occupational use supports S; principles or theory support a complete K phrase. Familiarity alone does not determine type.\gsource{Q,P} & \gcontext{使用}\gspan{S}{SQL}\ggloss{Use SQL}\par\smallskip\gcontext{了解}\gspan{K}{SQL原理}\ggloss{Understand the principles of SQL} \\[5pt]
G12 & Occupational functions are distinguished from general cognition, communication, and coordination. A duty clause may contain T; usefulness across industries alone does not establish T.\gsource{E,P} & \gspan{S}{编码能力}\ggloss{Coding skills}\par\smallskip\gspan{T}{沟通能力}\ggloss{Communication skills}\par\smallskip\gspan{T}{协调外部资源}\ggloss{Coordinate external resources} \\[5pt]
G13 & Scope, boundaries, and type are resolved separately, without a fixed priority among labels. If only the type remains uncertain, a mention may receive provisional S with alternatives recorded. Such unresolved records are excluded from final training targets.\gsource{I,P} & Heading alone: \zh{能力提升}\ggloss{Developing capabilities}\par\smallskip The heading alone leaves scope unresolved; it does not establish an empty annotation. \\[5pt]
G14 & Offsets count original Unicode code points, beginning at zero and excluding the end position. The indexed substring must reproduce the mention exactly; segmentation is a derived view. It does not replace the fixed human-reference offsets.\gsource{P} & \texttt{text[start:end] == span}\par\smallskip Keep the original text, internal spaces, and offsets consistent. \\[5pt]
G15 & Alternatives and overlaps are recorded separately from the final flat layer. An unresolved record is excluded in full from training, with confirmed empty, missing, and unresolved annotations distinguished.\gsource{P} & \zh{五险一金}\ggloss{Five types of social insurance and a housing fund}\par $\rightarrow$ confirmed \texttt{[]};\par\smallskip\zh{能力提升}\ggloss{Developing capabilities}\par Without context $\rightarrow$ unresolved.\par\smallskip Known partial spans do not make an unresolved record training-ready. \\
\bottomrule
\end{longtable}
\endgroup

\paragraph{Interpreting coordinated capabilities.}
G6--G7 test whether each selected phrase expresses a requirement in context. Items need not repeat an implied capability noun. For example, \zh{协调} (\emph{coordination}) can express a capability independently, whereas \zh{曝光} (\emph{exposure}) alone does not express the optimization action in \zh{优化曝光与转化率} (\emph{optimize exposure and the conversion rate}). The figure preserves the human-reference boundaries used for scoring; the examples below illustrate the rules summarized here. Their respective handbook versions are identified in Section~A.4.

\subsection{A.3 Shared experience: splitting without changing the requirement}
The V4.2.14 amendment (R23) clarifies a shared-experience construction.\gsource{P} Each item remains identifiable in context; exclude the shared experience qualifier and do not add a covering span. The words for internet-company operations denote an occupational activity in this sentence; a bare requirement for experience in the internet industry would instead use the project K proxy rule. The presence of an industry name alone does not determine type.

\smallskip\noindent
\begin{minipage}{\linewidth}\small\RaggedRight
\textbf{Chinese.}\enspace\gcontext{4、三年以上大中型}\gspan{S}{互联网公司经营}\gcontext{或}\gspan{S}{BI}\gcontext{相关工作经验，}\gspan{S}{互联网公司广告运营分析}\gcontext{优先；}
\par\smallskip
\textbf{English.}\enspace At least three years of relevant work experience in \gspan{S}{internet company operations} or \gspan{S}{BI} at medium-sized or large companies; experience in \gspan{S}{advertising operations analysis at internet companies} preferred.
\par\smallskip
Matching highlights align the three S spans across languages. Chinese offsets: $[9,16)$, $[17,19)$, and $[26,37)$; the original item number is retained. BI is S in this context, not by a universal dictionary rule. In contrast, G7 retains the shared action in \zh{优化曝光与转化率} (\emph{optimize exposure and the conversion rate}); splitting must not leave bare outcomes. The initial reviewed LLM-annotation layer retains the earlier long span.
\end{minipage}

\subsection{A.4 Handbook Versions Used in This Study}
\label{sec:guide-version-history}
Table~\ref{tab:handbook-stage-records} links annotation stages to their documented handbook versions. Label metadata and operating instructions provide different evidence of version use; individual silver training and development records have no handbook-version field.

\begin{table}[!htbp]
\caption{Handbook versions associated with the reported annotation stages.}
\label{tab:handbook-stage-records}
\centering\small
\setlength{\tabcolsep}{4pt}
\begin{tabularx}{\linewidth}{@{}P{.22\linewidth}P{.23\linewidth}L@{}}
\toprule
Annotation stage & Handbook record & Interpretation \\
\midrule
Human reference & v4.2.10; v4.2.11 & All 150 records carry v4.2.10 metadata; reference-set documentation identifies v4.2.11 as an operating handbook. \\
Initial silver review & v4.2.10 rev1; v4.2.11 & The initial 100-sentence prompt uses v4.2.10 rev1; later coding and adjudication records identify v4.2.11. \\
Bulk silver generation & v4.2.12 rev2 & Basis of the archived prompt for the remaining 2,351 original silver records. \\
Blinded-study preparation and coding & v4.2.14 series & Training materials are preserved. Exact files used for formal coding and individual self-review have not been established. \\
\bottomrule
\end{tabularx}
\end{table}

The \href{https://github.com/AlfredJamesLi/chinese-skillspan-benchmark/blob/main/reproduction/silver_prompt_api_v1/archive/silver_rev2_original.txt}{bulk annotation prompt} and \href{https://github.com/AlfredJamesLi/chinese-skillspan-benchmark/blob/main/notes/handbooks/README.md}{handbook index} provide the study materials and dated revisions. Several revisions share the v4.2.14 number and are distinguished by edition date and file. Later reader editions clarify rules for new annotation without replacing preserved study labels.

The shared-experience example in Section~A.3 illustrates a later boundary clarification; its original training annotation is retained. Figure~\ref{fig:illustrative-annotations} reproduces the human-reference labels, whereas Table~\ref{tab:guide-types} and Tables~\ref{tab:guide-boundaries}--\ref{tab:guide-scope} illustrate the v4.2.14 rules summarized here. Batch records, version history, and the corresponding example are documented in the \href{https://github.com/AlfredJamesLi/chinese-skillspan-benchmark/blob/main/reproduction/annotation_documentation_20261007/README.md}{annotation documentation}.

Archived predictions can be rescored with their reference and the fixed scorer. New model outputs or annotations may vary with handbook text, model versions, or service settings.
